% Supplementary material: how to read a claim register.
% Body only. The preprint build (\preprintmode) includes it as an appendix of
% manuscript.tex; supplement-standalone.tex builds it as the separate
% supporting-information file Learned Publishing takes. Not counted against the
% manuscript's word budget (\InputIfFileExists). Its claims are not yet registered.

\section{How to read a claim register}
\label{app:reading}

The claim register builds each row on Toulmin's layout of arguments.  This
supplement introduces the layout for readers who have not met it, and reads
one row of this paper's register with it.

\subsection{Toulmin's layout of an argument}

Toulmin modelled his layout of arguments on jurisprudence, calling logic
``generalised jurisprudence'' \citep[p.~7]{toulmin2003}.  He takes apart an
everyday argument \citep[pp.~92--94]{toulmin2003}:

\begin{center}
\small
\begin{tabular}{@{}c@{\quad}c@{\quad}c@{}}
\textit{Harry was born in Bermuda} & $\longrightarrow$ & \textit{so, presumably, Harry is a British subject} \\
(data) & & (qualifier, claim) \\[6pt]
\multicolumn{3}{c}{the step holds \emph{since} \textit{a man born in Bermuda will generally be a British subject}\ (warrant)} \\[2pt]
\multicolumn{3}{c}{\emph{unless} \textit{both his parents were aliens, or he has become a naturalised American}\ (rebuttal)} \\
\end{tabular}
\end{center}

\begin{itemize}
\item The \textbf{claim} is what the arguer asks us to accept.
\item The \textbf{data} are ``the facts we appeal to as a foundation for
  the claim'' \citep[p.~90]{toulmin2003}.  The register calls them
  \emph{grounds}, because they can also be other claims, a declared
  starting point or the author's own work.
\item The \textbf{warrant} is one of the ``general, hypothetical statements,
  which can act as bridges'' from data to claim \citep[p.~91]{toulmin2003}:
  the general rule that licenses the step.  Its task is ``simply to register
  explicitly the legitimacy of the step involved'', and ``data are appealed
  to explicitly, warrants implicitly'' \citep[p.~92]{toulmin2003}.  In a
  paper the warrant is often unstated, and that is where a weak step hides.
\item The \textbf{qualifier} indicates ``the strength conferred by the warrant
  on this step'' \citep[p.~94]{toulmin2003}: here, ``presumably''.
\item The \textbf{rebuttal} names ``circumstances in which the general
  authority of the warrant would have to be set aside''
  \citep[p.~94]{toulmin2003}.
\item The \textbf{backing} is what supports the warrant itself, such as the
  statutes behind the rule about Bermuda.  The register records it only when
  readers may not accept the warrant.
\end{itemize}

Toulmin compares the distinction between data and warrants to ``the
distinction drawn in the law-courts between questions of fact and questions
of law'' \citep[p.~92]{toulmin2003}: data are the facts of the case, and the
warrant is the law applied to them.  He later wrote that he had not had in
mind an analytical model like the one that came to be called ``the Toulmin
model''; his concern ``was with twentieth-century epistemology, not informal
logic'' \citep[p.~vii]{toulmin2003}.

\subsection{One row of this paper's register}

Table~\ref{tab:s1-1} reads row S1-1 with the layout.

\begin{table}[htbp]
\centering
\caption{Register row S1-1 read with Toulmin's layout.}
\label{tab:s1-1}
\footnotesize
\begin{tabular}{@{}p{2.6cm}p{9.4cm}p{2.0cm}@{}}
\toprule
\textbf{Register column} & \textbf{Content} & \textbf{Toulmin} \\
\midrule
Statement & LLMs tested in 2023--24 (ChatGPT on GPT-3.5 and GPT-4) invented plausible-sounding references and DOIs at substantial rates. & Claim \\
Grounds & Two peer-reviewed sources, each reporting its own study \citep{mugaanyi2024,walters2023}. & Data \\
Warrant & None needed: each source reports the claim's own point. & Warrant \\
Confidence tier & Supported: two peer-reviewed sources agree, so the prose may say ``indicates'', not ``shows''. & Qualifier \\
Rebuttal & As of 2023--24; one study found fewer fabrications from the newer model, so later models may fabricate less. & Rebuttal \\
\bottomrule
\end{tabular}
\end{table}

The first version of this row said that LLMs ``invent plausible-sounding
papers'': present tense, all models.  Nothing in the grounds licensed the
step from two studies of ChatGPT in 2023--24 to every current model.  In
Toulmin's terms the row had an unstated warrant, and the warrant was false.
A reader caught it, not a citation check.

\subsection{Grounds that are not sources}

A \textbf{declared starting point} is an assumption, definition or value that
\emph{this paper} adopts without arguing for it, for example D9: a ``citation check'' is a check that a cited work exists and says what is attributed to it.  It is not a
claim taken from a source.  A point that a source took from another source
is still a source item, recorded as relayed (the source check, Section~4 of
the paper).

The full register, the glossary of terms and the tools that check them are
in the public repository \citep{veen2026repo}.
